














\section{Introduction}
\label{sec:introduction}

Tissue function depends not only on local cellular composition, but also on how cells and their neighborhoods are organized within larger anatomical structures~\citep{schurch2020coordinated}. Spatial transcriptomics measures molecular profiles together with spatial position, making it possible to study these levels of tissue organization within the same specimen~\citep{chen2022mosta,pan2026hesta}. Importantly, tissue characteristics depend on the spatial scale at which they are considered: local neighborhoods capture cellular composition and interactions, whereas broader spatial contexts reveal regional organization and anatomical structure. Relating these scales is therefore essential for understanding how local molecular environments are embedded within larger tissue organization.

Spatial representation methods typically characterize tissue context within a predefined neighborhood or spatial graph~\citep{long2023graphst,singhal2024banksy}. Multiscale methods extend this view by learning representations at several spatial scales or combining information across them~\citep{yuan2024mender,zhao2025stofm,lin2024stalign}. However, these approaches largely provide representations at a discrete set of predefined scales and treat them as parallel sources of information to be combined. Tissue organization can change systematically with scale. \textbf{Local structures may persist as broader spatial context introduces new regional characteristics and reorganizes their spatial relationships.} Capturing such scale-dependent organization therefore requires modeling not only representations at individual scales, but also how they transform across scales.

Existing pretraining objectives provide limited supervision for learning these transformations. Single-cell pretraining commonly predicts masked genes or expression values from molecular context within a cell~\citep{theodoris2023geneformer,cui2024scgpt}, while spatial predictive objectives incorporate neighboring cells to recover molecular information or latent cell states~\citep{zhao2025stofm,liu2026cellworld}. These objectives primarily learn what can be inferred within a given context, rather than how representations vary with spatial scale. Organizing scale-specific representations within a shared latent space makes their common and scale-dependent characteristics directly comparable, while modeling transformations between them extends representation learning beyond a few predefined scales toward continuously varying spatial contexts. This enables a more fine-grained characterization of how tissue context changes with spatial scale.


We introduce ZoomST, a self-supervised pretraining framework for spatial transcriptomics that learns molecular--spatial dual-view representations of tissue context across spatial scales. We design a soft distributional equivariance constraint to learn structured transformations between scale-specific representations within a shared latent space while preserving scale-dependent tissue characteristics. This yields a structured cross-scale representation space that extends beyond a discrete set of embeddings at predefined spatial granularities, enabling tissue context to be characterized continuously across spatial granularities.

We evaluate both the representations and their transformations in developmental spatial transcriptomic atlases. Our main contributions are:

\begingroup
\setlength{\leftmargini}{1.3em}
\begin{itemize}
\item \textbf{Molecular--spatial dual-view pretraining for tissue representation.}
We develop a self-supervised pretraining framework that jointly learns molecular and spatial views of tissue context across multiple spatial scales, and provides scale-specific representations in a shared latent space.

\item \textbf{Soft distributional equivariance for structured cross-scale learning.}
We introduce a soft distributional equivariance constraint that imposes cross-scale consistency within a shared latent space while preserving scale-dependent tissue characteristics. This enables ZoomST to learn structured transformations between representations rather than treating different scales as independent embeddings.

\item \textbf{Continuous cross-scale transformation and scale sensitivity.}
The learned transformations interpolate tissue representations at unseen intermediate spatial scales despite supervision at only a discrete set of scales. They further provide a means to characterize how strongly tissue representations respond to changes in spatial context, revealing spatial variation in scale sensitivity.

\item \textbf{Generalization across developmental spatial transcriptomic atlases.}
We validate ZoomST on human and mouse developmental atlases. Frozen representations generalize to held-out tissue sections and support downstream tasks including tissue mapping and developmental organization, outperforming the evaluated pretrained baselines on tissue mapping while remaining competitive with target-fitted methods.
\end{itemize}
\endgroup
